% \newpage
\section{Preliminaries}
\label{sec: mcopi}

% Monte Carlo optimistic policy iteration (MC-O-PI) is a reinforcement learning algorithm that can learn the optimal policy of a finite Markov decision process (MDP) without knowing the underlying parameters of that MDP.
% This section briefly introduces MDPs, then presents the stochastic difference inclusion that underlies MC-O-PI.

\subsection{Finite Markov decision processes}


% Reinforcement learning studies the interaction between an agent and an environment.
% ...
% this can be modelled using finite MDP
% ...
% At each time step \(t\), the agent observes a state \(s_t \in \mathcal S\) and selects an action \(a_t \in \mathcal A(s_t)\) according to a policy \(\pi(a_t | s_t)\).
% The environment then produces a reward \(r_t\) and a new state \(s_{t+1}\) according to its transition dynamics $p(s_{t+1}, r_t | s_t, a_t)$.
% % 
% The action value of a policy \(\pi\) for any pair $(s,a)$ is the expected discounted return obtained by taking action \(a\) in state \(s\), and then following policy \(\pi\):
% \[
% q_\pi(s,a)
% =
% \mathbb E \left[
% \sum_{t=0}^{\infty} \gamma^t r_t
% \;\middle|\;
% s_0 = s,\
% a_0 = a,\
% a_t \sim \pi(\cdot \mid s_t),\
% (s_{t+1},r_t) \sim p(\cdot,\cdot|s_t,a_t), \
% \forall\, t\ge 1
% \right].
% \]
% Under standard finite-MDP assumptions,
% there exists an optimal policy $\pi_*$ that satisfies $q_{\pi_*}(s,a) = \max_\pi q_\pi(s,a)$ for all $(s,a)$.
% Moreover, $\pi_*$ can be chosen greedily with respect to $q_{\pi_*}$, meaning
% % $q_{\pi_*}(s,\pi_*(s)) = \max_a q_{\pi_*}(s,a)$ for all $s$.
% $\pi_*(s) \in \arg\max_a q_{\pi_*}(s,a)$ for all $s$.


A Markov decision process evolves over discrete steps.
% with finite state space $\sset$ and finite action space $\aset(s)$ for each state $s \in \sset$.
At time \(t\), the agent observes a state \(s_t \in \mathcal S\) and selects an action \(a_t \in \mathcal A(s_t)\) according to a policy \(\pi(a_t | s_t)\).
The environment then produces a reward \(r_t\) and transitions to a new state \(s_{t+1}\) according to its dynamics $p(s_{t+1}, r_t | s_t, a_t)$.
% Taking action $a$ in state $s$ yields a finite reward $r(s,a) \in \real$ 
% and transitions to a new state $s'$ with probability $p(s' \mid s, a)$.
% The discount factor is $\discount \in [0,1]$.
If a policy $\pol$ is deterministic in state $s$, then $\pol(s)$ directly refers to the action $a$ where $\pol(a | s) = 1$.
We denote the set of deterministic policies by $\polset$.

% A stationary policy is a map $\pol(\cdot \mid s)$ that assigns a distribution over $\aset(s)$ for each state $s$.
% As is common in the literature, we refer to that distribution using the notation $\pol(s)$.
% We also insert the distribution $\pol(s)$ in functions with domain $\aset(s)$.
% For instance, the expected immediate reward when following policy $\pol$ in state $s$ is
% \begin{align*}
%     r(s, \pol(s)) \equiv \sum_{a \in \aset(s)} \pol(a \mid s) r(s, a) .
% \end{align*}
% If a policy $\pol$ is deterministic in state $s$, then $\pol(s)$ directly refers to the action $a$ where $\pol(a \mid s) = 1$.
% The set of deterministic policies is denoted by $\polset$.



% \begin{align}
%     \polset \coloneqq 
%     \Big\{ \pol : \aset \times \sset \to \{ 0,1 \}
%     \ \Big| \
%     \forall s \in \sset : \sum_{a \in \aset(s)} \pol(a \mid s) = 1 
%     \Big\} .
% \end{align}

% Following a policy $\pi$ generates an episode, namely a sequence $(s_t, a_t, r_t)_{0 \leq t < T}$ such that for each $t < T$
% \begin{align*}
%     a_t \sim \pol(\cdot \mid s_t) ,
%     \quad\quad
%     r_t = r(s_t, a_t) ,
%     \quad\quad
%     \ s_{t+1} \sim \transit(\cdot \mid s_t, a_t) .
% \end{align*}
% The return of the episode for a state-action pair $(s_t, a_t)$ is
% \begin{align*}
%     r_t + \discount r_{t+1} + \discount^2 r_{t+2} + \dots + \discount^{T-1-t} r_{T-1} .
% \end{align*}
% These episodes induce the standard value functions, as
% the state-value function $v_\pol : \sset \to \real$ measures the expected return
% when starting in state $s$ and following policy $\pol$:
% \begin{align*} 
%     v_\pol(s) 
%     \coloneqq 
%     \E_{\transit} \big[ r(s_0, \pol) + \discount r(s_1, \pol) + \discount^2 r(s_2, \pol) + \dots \mid s_0 = s \big] ,
% \end{align*}
For any policy $\pi$,
the state-value function $v_\pol : \sset \to \real$ 
measures the expected discounted return when following policy $\pol$ starting in state $s$:
\begin{align*} 
    v_\pol(s)
    :=
    % \E_{p} \big[ r(s_0, a_0) + \discount r(s_1, \pol)
    % \\
    % &\hspace{4.5cm} + \discount^2 r(s_2, \pol) + \dots \mid s_0 = s, a_0 = a \big] .
    \mathbb E_{p, \pi} \left[ \sum_{t=0}^{\infty} \gamma^t \, r_t  \ \bigg| \ s_0 = s \right] .
\end{align*}
The action-value function $q_\pol : \sset \times \aset \to \real$ measures the expected discounted return when taking action $a$ in state $s$ and then following policy $\pol$:
\begin{align*} 
    q_\pol(s,a)
    :=
    % \E_{p} \big[ r(s_0, a_0) + \discount r(s_1, \pol)
    % \\
    % &\hspace{4.5cm} + \discount^2 r(s_2, \pol) + \dots \mid s_0 = s, a_0 = a \big] .
    \mathbb E_{p, \pi} \left[ \sum_{t=0}^{\infty} \gamma^t \, r_t  \ \bigg| \ s_0 = s, a_0 = a \right] .
\end{align*}

A policy $\pi$ is at least as good as another policy $\pi'$ if $v_{\pi} \ge v_{\pi'}$ or, equivalently, $q_{\pi} \ge q_{\pi'}$.
We denote this relation by $\pi \succeq \pi'$.
If, additionally, there exists a state where $v_{\pi}(s) > v_{\pi'}(s)$, then $\pi$ is better than $\pi'$, denoted $\pi \succ \pi'$.


% For a policy \(\pi\), the action-value function $q_\pi$ assigns to each pair $(s,a)$ the expected discounted return obtained by taking action \(a\) in state \(s\), and then following policy \(\pi\):
% \[
% q_\pi(s,a)
% =
% \mathbb E_{\pi} \left[
% \sum_{t=0}^{\infty} \gamma^t r_t
% \;\middle|\;
% s_0 = s,\
% a_0 = a
% % a_t \sim \pi(\cdot \mid s_t),\
% % (s_{t+1},r_t) \sim p(\cdot,\cdot|s_t,a_t), \
% % \forall\, t\ge 1
% \right].
% \]

% Under standard finite-MDP assumptions,
% there exists an optimal policy $\pi_*$ that satisfies $q_{\pi_*}(s,a) = \max_\pi q_\pi(s,a)$ for all $(s,a)$.
% Moreover, $\pi_*$ may be chosen to be deterministic and greedy with respect to $q_{\pi_*}$, meaning
% $\pi_*(s) \in \arg\max_{a \in \mathcal A(s)} q_{\pi_*}(s,a)$ for all $s$.


This paper considers discounted or episodic MDPs,
meaning $\gamma \in [0,1)$
or episodes simulated with any stationary policy reach a terminal state almost surely.
The Policy Improvement Theorem then allows us to compare policies using only the action-value function of one of them~\citep{Howard1960DynamicProcesses}.
We rely on a strict version of that theorem.

% \label{thm: policy improvement theorem}\begin{theorem}
% [Policy Improvement Theorem]
% If policies $\pi$ and $\pi'$ satisfy
% \begin{align}
% % \label{eq: strict improvement theorem}
% q_\pi(s,\pi(s)) \ge q_\pi(s,\pi'(s))
% \end{align}
% for every state $s \in \sset$,
% then policy $\pi$ is at least as good as policy $\pi'$.
% If the inequality is strict in at least one state, 
% then policy $\pi$ better than policy $\pi'$.
% \end{theorem}

\begin{lemma}
\label{thm: strict PIT}
Given policies $\pi$ and $\pi'$, suppose that
\begin{align}
\label{eq: PIT condition}
q_\pi(s,\pi'(s)) \ge q_\pi(s,\pi(s))
% \qquad \forall s\in\mathcal S .
\end{align}
for every state $s\in\mathcal S$.
Then \(v_{\pi'} \ge v_\pi\). Moreover, if \(q_{\pi'}\neq q_\pi\), then the improvement is strict, meaning
there exists a state \(s\in\mathcal S\) such that
\[
v_{\pi'}(s)>v_\pi(s).
\]
\end{lemma}
\begin{proof}
By the Policy Improvement Theorem, inequality \eqref{eq: PIT condition} implies that $v_{\pi'} \ge v_\pi$.
Suppose that $v_{\pi'} = v_{\pi}$. Then for every pair $(s,a)$,
\begin{align*}
q_{\pi'}(s,a)
% = r(s,a) + \gamma \sum_{s' \in \mathcal S} p (s' \mid s,a) v_{\pi'}(s')
= \sum_{s'} \sum_{r} p (s', r \mid s,a) \big( r + \gamma v_{\pi'}(s') \big)
% = r(s,a) + \gamma \sum_{s' \in \mathcal S} p (s' \mid s,a) v_{\pi}(s')
= \sum_{s'} \sum_{r} p (s', r \mid s,a) \big( r + \gamma v_{\pi}(s') \big)
= q_{\pi}(s,a) .
\end{align*}
Thus, if $q_{\pi'} \neq q_\pi$, there exists at least one state where $v_{\pi'}(s) > v_{\pi}(s)$, and thus $\pi' \succ \pi$.
\end{proof}

Classical results guarantee that, for discounted or episodic MDPs, 
there exists an optimal policy $\pi_*$ that satisfies $q_{\pi_*}(s,a) = \max_\pi q_\pi(s,a)$ for all pairs $(s,a)$.
Moreover, $\pi_*$ may be chosen to be deterministic and greedy with respect to $q_{\pi_*}$, meaning
$\pi_*(s) \in \arg\max_{a \in \mathcal A(s)} q_{\pi_*}(s,a)$ for all $s$.
The converse is also true: if a deterministic policy is greedy with respect to its own action values, then it must be optimal.
This essential property allows us to verify whether a policy is optimal by inspecting its action values, rather than comparing it with all other policies.


% there exists at least one deterministic policy that is better than or equal to all other policies \citep{Bellman2010DynamicProgramming}.
% We denote the set of all such optimal policies by $\mathcal P_*$. By definition, all optimal policies share the same value functions, denoted $v_{\best{\pi}}$ and $q_{\best{\pi}}$.

% Let $\pi_* \in \mathcal P_*$ be a deterministic optimal policy.
% The Policy Improvement Theorem implies that $\pi_*$ is greedy with respect to its own action values:
% \begin{align}
%     q_{\pi_*}(s, \pi_*(s)) = \max_{a \in \mathcal A(s)} q_{\pi_*}(s,a)
%     \qquad \forall \, s \in \mathcal S .
% \end{align}
% The converse is also true: if a deterministic policy is greedy with respect to its own action values, then it must be optimal.
% This property is essential because it allows us to verify whether a policy is optimal by inspecting its action values, rather than comparing it with all other policies.


% As a result, the optimal action values satisfy
% \begin{align*}
%     q_{\best{\pi}}(s,a) = \max_{\pi \in \polset} q_{\pi}(s,a),
%     \qquad \forall s \in \sset, \forall a \in \aset(s) .
% \end{align*}

% The optimal policy $\best{\pi}$ thus maximises the value functions:
% % for all state–action pairs simultaneously:
% \begin{align*}
%     &v_{\best{\pi}}(s) = \max_{\pi \in \polset} v_{\pi}(s)
%     % \\
%     &q_{\best{\pi}}(s,a) = \max_{\pi \in \polset} q_{\pi}(s,a),
%     \qquad \forall s \in \sset, \forall a \in \aset(s) .
% \end{align*}



% for every state $s$ and action $a$.
% The objective of MC-O-PI is to find an optimal policy $\best{\pi}$. 
% This policy maximises the value functions for all states (equivalently, all state–action pairs):
% \begin{align}
%     &v_{\best{\pi}}(s) = \max_{\pi} v_{\pi}(s)
%     \\
%     &q_{\best{\pi}}(s,a) = \max_{\pi} q_{\pi}(s,a)
% \end{align}
% for every state $s$ and action $a$.
% % Note that the maximum is attained because the set of policies is bounded.


% A deterministic policy is greedy with respect to a function $q$ if it selects only maximisers of $q$ in each state:
% \begin{align}
%     \pi(s) \in \argmax_{a \in \aset(s)} q(s,a)
% \end{align}


% Since the optimal policy is deterministic and the set of deterministic policies is finite, one could find it by enumerating all deterministic policies,
% solving the linear system $q = T_\pi q$ for $q$, and verifying the optimality condition $q_\pi = T q_\pi$. However, this approach is inefficient as it ignores previously gathered information to guide the search. Moreover, it is infeasible when the parameters of the MDP are unknown. In such situations, practical, model-free learning algorithms are necessary.


% Consequently, every policy that is greedy with respect to its own action values is optimal.

% \begin{definition}[Bellman operators]
% \label{def: bellman operators}
% The transition operators $\T_\pol, \T : \qset \to \qset$ are defined as
% \begin{align*}
%     (\T_{\pol} q)(s,a) &\coloneqq \sum_{s_+ \in \sset} \transit(s_+ \mid s,a) q(s_+,\pol) ,
%     \\
%     (\T q)(s,a) &\coloneqq \sum_{s_+ \in \sset} \transit(s_+ \mid s,a) \max_{a_+ \in \aset(s_+)} q(s_+, a_+) .
% \end{align*}
% The Bellman operators $\B_\pol, \B : \qset \to \qset$ are defined as
% \begin{align*}
%     (\B_\pol q) (s,a) &\coloneqq \reward(s,a) + \discount \sum_{s_+ \in \sset} \transit(s_+ \mid s,a) q(s_+, \pol) ,
%     \\
%     (\B q)(s,a) &\coloneqq \reward(s,a) + \discount \sum_{s_+ \in \sset} \transit(s_+ \mid s,a) \max_{a_+ \in \aset(s_+)} q(s_+,a_+) .
% \end{align*}
% \end{definition}

% In this paper we assume the discount factor $\discount<1$ so that for each deterministic policy $\pi$ the policy Bellman operator $\B^\pi$ is monotone with nonnegative linear part, $q_\pi$ is its unique fixed point, and whenever $\pi\in P(q)$ (i.e.\ $\pi$ is greedy w.r.t.\ $q$) one has $\B q = \B^\pi q$.


%%%%%%%%%%%%%%%%%%%%%%%%%%%%%%%%%%%%%%%%%%%%%%%%%%%%%%%%%%%%%%%%%%%%%%%%%%
%%%%%%%%%%%%%%%%%%%%%%%%%%%%%%%%%%%%%%%%%%%%%%%%%%%%%%%%%%%%%%%%%%%%%%%%%%
\subsection{Monte Carlo optimistic policy iteration}

In many real-world applications, the transition dynamics and reward functions of the underlying MDP are unknown. Learning the optimal policy therefore requires interacting with the system to estimate expected returns from accumulated experience.
One fundamental framework for this is \textit{policy iteration}.
This approach tracks an action-value estimate $q$ and a policy $\pol$ and refines them by alternately evaluating the current policy and improving it.
This cycle is intended to converge to a policy that is greedy with respect to its own action values, and therefore optimal.


Monte Carlo optimistic policy iteration implements the policy-iteration framework by using 
\textit{Monte Carlo policy evaluation}, which simulates complete episodes of the MDP to update $q$ towards the action values of $\pol$,
and \textit{optimistic policy improvement}, which updates $\pol$ to obtain higher returns as anticipated by $q$.


Specifically, consider the function space $\qset \coloneqq \{ \mathcal S \times \mathcal A \to \mathbb R \}$.
% and let $k$ be an iteration index.
MC-O-PI generates a sequence $(q_k)_{k \ge 0}$ in $\qset$ as follows.
At the beginning of iteration $k$, the algorithm selects a deterministic policy $\pi_k$ that is greedy with respect to $q_k$:
\begin{align}
\label{eq: greedy policy definition}
    \pol_k(s) \in \arg\max_{a \in \aset(s)} q_k(s,a),
    \qquad \forall \, s \in \sset .
\end{align}
The algorithm then samples an initial state-action pair $(s_0, a_0)$ according to a sampling distribution $\start_k$,
simulates an episode of the MDP following policy $\pol_k$, and observes the discounted return for each state-action pair in the episode.
% Since the expected return for a given pair is $q_{\pi_k}(s,a)$, t
The observed return equals $q_{\pi_k}(s,a) + v_k(s,a)$, where $v_k(s,a)$ is a stochastic perturbation.
Finally, the estimate $q_k$ is updated at pairs appearing in the episode
according to
\begin{align}
\label{eq: MC eval update}
    q_{k+1}(s, a) 
    &= q_k(s,a) +  \lrate_k( q_{\pol_k}(s, a) + v_k(s, a) - q_k(s,a) ) ,
    % \\
    % &= q_k(s,a) + \lrate_k ( q_{\pol_k}(s, a) + v_k(s, a) - q_k(s,a)) .
\end{align}
where $(\lrate_k)_{k \ge 0}$ is a scalar learning rate sequence that satisfies the Robbins-Monro conditions:
\begin{align}
 \label{eq: robbins monro conditions}
    % 0 < \alpha_k \le 1 ,
    % \qquad
    \sum_{k =0}^\infty \alpha_k = \infty , 
    \qquad
    \sum_{k =0}^\infty \alpha_k^2 < \infty .
\end{align}
% $\sum_{k \ge 0} \lrate_k = \infty$ and $\sum_{k \ge 0} \lrate_k^2 < \infty$.


Let $\mathcal F_k$ represent the available information after $q_k$ is computed,
but before
selecting the greedy policy, 
sampling the initial state-action pair 
and simulating the episode.
% information available before simulating the $k^{\text{th}}$ episode and drawing the $k^{\text{th}}$ update mask, namely
% \mytodo{why only know $\pi_i$ up to k-1? F should contain $\pi_k$}
% \mytodo{showd $f_k$ and $\sigma_k$ also be in the filtration}
% \begin{align}
% \label{eq: filtration_def}
% % \mathcal F_k
% % \coloneqq
% % \sigma
% \Big\{
% q_0,
% \big(\alpha_i, \sigma_i, f_i, q_i,\pi_i,(s^{(i)}_t,a^{(i)}_t,r^{(i)}_{t})_{0\le t<T_i} \big)_{0\le i\le k-1}
% \Big\}.
% \end{align}

\subsubsection{Greedy policy}

Define the set-valued function $\grpol : \qset \rightrightarrows \polset$ that maps any action-value function $q$ to the set of deterministic policies that are greedy with respect to $q$:
\begin{align}
\label{eq: greedy pol}
    \grpol(q) \coloneqq \big\{ \pol \in \polset : q(s, \pol(s)) = \max_{a \in \aset(s)} q(s,a), \ \forall \, s \in \sset \big\} .
\end{align}
At every step, the policy $\pi_k$ thus satisfies 
\begin{align}
\label{eq: greedy policy def}
    \pi_k \in \grpol(q_k) .
\end{align}
If several policies are greedy with respect to $q_k$, we assume that ties are broken randomly and independently across states.
More precisely, we assume that
there exists a constant \(\beta>0\) such that
% for every time step \(k \ge 0\), every state \(s\in S\), and every action $a \in \arg\max_{A(s)} q_k(s,\cdot)$,
\begin{equation}
\label{eq:greedy-pol-tie-breaking}
\mathbb P(\pi_k(s)=a\mid \mathcal F_k)\ge \beta,
    \qquad 
    \forall \, k \ge 0, \
    \forall \, s \in \sset, \
    \forall \, a \in \arg\max_{\mathcal A(s)} q_k(s,\cdot) .
\end{equation}
% For example, under uniform tie-breaking,
% \[
% \mathbb P(\pi_k(s)=a\mid \mathcal F_k)
% =
% \frac{1}{|\arg\max_{a'\in A(s)} q_k(s,a')|}
% \ge \frac{1}{|A(s)|},
% \]
% so \eqref{eq:greedy-pol-tie-breaking} holds with
% \[
% \beta=\min_{s\in S}\frac{1}{|A(s)|}.
% \]
This assumption is mostly technical as ties rarely occur in practice, except if the algorithm is initialised in one.


\subsubsection{Update function}

Not all state-action pairs are necessarily updated after each episode.
Given an episode $\big( (s_t, a_t) \big)_{t \ge 0}$, there are three common update functions to decide which pairs to modify.
The \textit{initial-visit} method only updates the first state-action pair $(s_0, a_0)$.
The \textit{first-visit} method updates each state-action pair $(s_t, a_t)$ provided it has not appeared in the sequence $\big((s_{t'}, a_{t'})\big)_{t' < t}$.
The \textit{every-visit} method updates all state-action pairs of an episode using every observed return from each pair.

We only consider initial- and first-visit updates to continue
because every-visit updates can produce biased estimates of the action values~\cite[Theorem~7]{Singh1996ReinforcementTraces}.
In other words, under initial- or first-visit updates the stochastic perturbations $(v_k)_{k \ge 0}$ are unbiased.
Their variance is also bounded.
Since both properties hold uniformly over the policies, it follows that for every time $k \ge 0$ and every pair $(s,a)$ visited in the $k^\nth$ episode,
\begin{align}
&\mathbb E[v_k(s,a) \mid \mathcal F_k] = 0 ,
\\
&\mathbb E[v_k^2(s,a) \mid \mathcal F_k] \le c_v ,
\end{align}
for some constant $c_v \in [0,\infty)$ that depends on the parameters of the MDP, but is independent of the policy $\pi_k$.
% \begin{align*}
% \mathbb E[v_k(s,a) \mid \mathcal F_k, \pi_k] = 0 .
% \end{align*}
% Furthermore, there exists a constant $c_v \in [0,\infty)$ such that
% \begin{align*}
% \mathbb E[v_k^2(s,a) \mid \mathcal F_k, \pi_k] \le c_v .
% \end{align*}
% Since both properties hold uniformly over the policies, we even obtain
% \begin{align}
% &\mathbb E[v_k(s,a) \mid \mathcal F_k] = 0 ,
% \\
% &\mathbb E[v_k^2(s,a) \mid \mathcal F_k] \le c_v .
% \end{align}
Under initial- or first-visit updates, MC-O-PI thus equates to
\begin{align}
\label{eq: update q with chi}
    q_{k+1}(s,a) 
    =
    q_k(s,a) + \lrate_k u_k(s,a) (q_{\pol_k}(s, a) + v_k(s, a) - q_k(s,a))   
    % \begin{cases}
    %     \ (1-\lrate_k)q_k(s,a) + \lrate_k (q_{\pol_k}(s, a) + v_k(s, a)) &\text{if} \ u_k(s,a)=1
    %     \\
    %     \ q(s,a) &\text{if} \ u_k(s,a)=0
    % \end{cases}
\end{align}
where $u_k(s,a)=1$ if $(s, a)$ is to be updated, and 0 otherwise.
The indicator function $u_k$ depends on the initial state-action pair of the $k^{\text{th}}$ episode, the policy used for simulating it and the chosen update function.
Considering the properties of $v_k$,
if we fixed the policy $\pi_k = \pi$ for all $k$, solutions of \eqref{eq: update q with chi} would converge almost surely to $q_{\pi}$ as long as all state-action pairs are updated infinitely often~\citep[Proposition~4.1]{Bertsekas1996Neuro-DynamicProgramming}.



\subsubsection{Initial sampling distribution}

To ensure that all state-action pairs are updated sufficiently often,
we commonly impose \textit{exploring starts} on the sampling distribution of the initial state-action pair of each episode. This condition requires that every state-action pair has a non-zero probability of being selected as the start of an episode:
% In other words, for every index $k>0$, state $s \in \sset$ and action $a \in \aset(s)$, the sampling distribution $\start_k$ of the initial state-action pair satisfies
\begin{align*}
    \start_k(s,a) > 0,
    \qquad 
    \forall \, k \ge 0, \
    \forall \, s \in \sset, \
    \forall \, a \in \aset(s) .
\end{align*}
However,
we can effectively nullify this condition using quickly decaying sampling distributions.
For example, if we choose $\sigma_k$ such that $\start_k(s,a) \le 1/(k+2)^2$ for some pair $(s,a)$, there is a $50\%$ chance that no episode ever starts in $(s,a)$. Under initial-visit updates, that pair would thus never be updated. To prevent such issues, we assume a stronger condition called \textit{strict exploring starts}:
there exists a fixed $\sigma_{\min} \in (0,\infty)$ such that
\begin{align}
\label{eq:minimal-sampling}    
\sigma_k(s,a) \ge \sigma_{\min},
    \qquad 
    \forall \, k \ge 0, \
    \forall \, s \in \sset, \
    \forall \, a \in \aset(s) .
\end{align}


\subsubsection{Update distribution}

The probability of updating each state-action pair after an episode depends on the distribution of the initial state-action pair, the policy used for simulating the episode and the update function.
We gather these probabilities in the update distribution $\mu_k$ defined as
\[
\mu_k := \mathbb E [ u_k \mid \mathcal F_k, \sigma_k, \pi_k, f_k] .
\]
% $: \sset \times \aset \to [0,1]$ 
% such that 
% $u_k = \mu_k + v'_k$, where $v'_k$ is zero-mean noise.

% \mytodo{$\mu_k$ is actually independent of $\mathcal F_k$ given the other three}

This distribution controls the exploration-exploitation trade-off
as it determines how frequently each action is updated.
Exploration favours actions whose current estimated value is most uncertain, for instance by biasing $\mu_k$ towards actions that received few updates in the past.
Exploitation, on the other hand, prioritises updating actions with the highest estimated value. It thus biases $\mu_k$ towards actions that are greedy with respect to $q_k$.


Using the update distribution, 
we can combine the stochastic terms $u_k$ and $v_k$ in equation \eqref{eq: update q with chi} into a single term $w_k$ defined as
\begin{align}
\label{eq: definition combined noise}
    w_k
    \coloneqq 
    u_k \, v_k + (u_k - \mu_k) ( q_{\pol_k} - q_k ) .
\end{align}
This term incorporates the perturbations caused by the episode simulation and the asynchronous updates of state-action pairs.
It simplifies equation \eqref{eq: update q with chi} to
\begin{align}
\label{eq: mcopi not yet sri}
    q_{k+1}
    =
    q_k + \lrate_k \big( \update_k ( q_{\pol_k} - q_k ) + w_k \big) .
\end{align}


% Thus the choice for the relative update probabilities of different actions controls the switching dynamics of MC-O-PI.
% This choice fundamentally reflects the exploration-exploitation trade-off.
% Exploitation prioritises updating actions with the highest current estimated value, which biases the update distributions towards the current greedy actions.
% Exploration on the other hand favours actions whose current estimated value is most uncertain, due to for instance fewer updates.
% Standard strategies provide concrete instances of this mechanism. For example, the epsilon-greedy approach with a decaying schedule starts with uniform update distributions,
% whereas the UCB-method initially favours actions that have received few updates.
% As the algorithm progresses, both methods progressively shift towards greedy update distributions.



% \textbf{Learning rates.}
% Unless stated otherwise, we assume that $\lrate_k \in \real_{> 0}$
% % In some examples however we use learning rates of the form $\lrate_k(s,a) = 1/n_k(s,a)$, where $n_k(s, a)$ denotes the number of times the estimate $q$ has been updated in the state-action pair $(s, a)$ during the first $k$ time steps. 
% % These learning rates are theoretically appealing because they ensure that $q_k(s,a)$ is the exact average of all approximations $\Tilde{q}_{\pol_k}(s,a)$ encountered up to step $k$.
% % However, they are rarely used in practice due to their higher memory requirements.
% % Regardless, the learning rates of stochastic approximation algorithms such as \eqref{eq: complete difference inclusion} must typically 
% and that the sequence $(\lrate_k)$ satisfies the Robbins--Monro conditions
% \begin{align}
%     \label{eq: Robbins--Monro first}
%     \sum_{k =0}^\infty \lrate_\kpz &= \infty , 
%     \\
%     \label{eq: robbins-monro second}
%     \sum_{k =0}^\infty \lrate_\kpz^2 &< \infty .
% \end{align}
% These conditions ensure that the learning rates are large enough to avoid premature convergence due to Zeno-like behaviour while gradually decreasing to zero so that noise cannot prevent convergence \citep{Robbins1951AMethod}.


\subsubsection{Difference inclusion}

% \begin{definition}[MC-O-PI difference ]
After defining the sequences of learning rates $(\lrate_k)_{k \ge 0}$, start distributions $(\start_k)_{k \ge 0}$ and update functions $(f_k)_{k \ge 0}$,
we can combine equations \eqref{eq: greedy policy def} and \eqref{eq: mcopi not yet sri} into the difference inclusion
\begin{multline}
\label{eq: mcopi sa}
    q_\kpo
    \in
    \Big\{ q_k + \alpha_k \big( \mu_k(q_{\pol_k} - q_k) + w_k \big) 
    :
    \\
    \pol_k \in \grpol(q_k),
    \update_k = \update(\start_k, \pol_k, f_k),
    w_k \sim w(\start_k, \pol_k, f_k)
    \Big\} .
\end{multline}
% \end{definition}
We assume the following conditions throughout this paper.

\begin{assumption}[Standing assumptions]
\label{asp: standing}
\phantom{Assume}
\begin{enumerate}[label=(\arabic*)]
    % \item \textbf{Discounted or episodic MDP.} 
    % % The Bellman operators $T_\pi$ and $T$ are monotonic and are $\gamma$-contractions in the $L_\infty$-norm.
    % Either the discount factor satisfies $\gamma \in [0, 1)$ or each episode of the MDP reaches a terminal state almost surely.
    
    \item \textbf{Robbins--Monro learning rates.} The learning rates satisfy condition \eqref{eq: robbins monro conditions}.
    % The learning rates $(\lrate_k)$ satisfy
    % \begin{align}
    % \label{eq: robbins monro conditions}
    %     % 0 < \alpha_k \le 1 ,
    %     % \qquad
    %     \sum_{k =0}^\infty \lrate_\kpz = \infty , 
    %     \qquad
    %     \sum_{k =0}^\infty \lrate_\kpz^2 < \infty .
    % \end{align}
    Since these conditions require that $\alpha_k \to 0$, we assume, without loss of generality, that $0 < \alpha_k \le 1$ for all $k \ge 0$.

    \item \textbf{Exploring policies.}
    The greedy policies satisfy condition \eqref{eq:greedy-pol-tie-breaking}.

    \item \textbf{Strict exploring starts.}
    The sampling distributions satisfy condition \eqref{eq:minimal-sampling}. 
    % \begin{align}
    % \label{eq: strict exploring starts}
    %     \sigma_k(s,a) \ge \sigma_{\min}
    %     \qquad 
    %     \forall \, k \ge 0, \
    %     \forall \, s \in \sset, \
    %     \forall \, a \in \aset(s) ,
    % \end{align}
    % for some constant $\sigma_{\min} \in (0,\infty)$.
    
    % \item \textbf{Unbiased episode return.} The update functions are initial- of first-visit updates.
\end{enumerate}
\end{assumption}


% where the set-valued function $\grpol : \qset \rightrightarrows \polset$ maps an action-value function $q$ to the set of deterministic policies that are greedy with respect to $q$:
% \begin{align}
% \label{eq: greedy pol}
%     \grpol(q) \coloneqq \big\{ \pol \in \polset : q(s, \pol(s)) = \max_{a \in \aset(s)} q(s,a), \ \forall s \in \sset \big\} .
% \end{align}

System \eqref{eq: mcopi sa} is a switched linear system whose active dynamics depend on which policy is greedy with respect to the current estimate $q_k$.
Thus, as long as the greedy policy does not change, the dynamics of \eqref{eq: mcopi sa} remain the same.
% which are essentially a stochastic approximation of the differential equation
% \[
% \dot q = \update_k ( q_{\pol_k} - q ) .
% \]
% 
We can delimit the subspaces of $\qset$ with constant dynamics using the set-valued function $\graval : \polset \rightrightarrows \qset$ defined as
\begin{align}
\label{eq: greedy action values}
    \graval(\pol) \coloneqq \big\{ q \in \qset : q(s, \pol(s)) = \max_{a \in \aset(s)} q(s,a), \ \forall \, s \in \sset \big\} .
\end{align}
This function is the inverse of $P$, defined in \eqref{eq: greedy pol}, as it maps any deterministic policy $\pol$ to the set of action-value functions for which $\pol$ is greedy.
Since $\graval(\pol)$ forms a convex subset of $\qset$, we refer to $\graval(\pol)$ as the \textit{region} of policy $\pol$ in $\qset$.


% For instance, consider the MDP in \autoref{fig: description region smallest MDP} with $\sset = \{s\}$ and $\aset(s) = \{a,a'\}$.
% Any point $q \in \qset$ can be represented by a point in $\real^2$ with coordinates $( q(s,a), q(s,a') )$.
% As shown in \autoref{fig: show regions}, the line $q(s,a) = q(s,a')$ bisects $\qset$.
% The region $\graval(\pol)$ is the half-space where $q(s,a) \ge q(s,a')$, while $\graval(\pol')$ is the half-space where $q(s,a') \ge q(s,a)$.

% Suppose the update distributions are uniform over all state-action pairs, and hence equivalent to a scalar in \eqref{eq: mcopi sa}.
% Following the ODE-interpretation,
% MC-O-PI solutions then approximate $\dot q = q_{\pi} - q$ over $\graval(\pi)$ and $\dot q = q_{\pi'} - q$ over $\graval(\pi')$, or in short
% \begin{align}
% \label{eq: switched ode for interpretation}
% \dot q \in \{ q_\pi - q : \pi \in \grpol(q) \}.
% \end{align}
% As the learning rate decreases, MC-O-PI follows the streamlines in \autoref{fig: explain regions with streamlines} more closely (arbitrary position of $q_\pi$ and $q_{\pi'}$).
% On the boundary, either systems are possible since the boundary belongs to both regions. However, the tie-breaking must satisfy \eqref{eq:greedy-pol-tie-breaking}.

% \input{figures/single-state-flow-explanation}



%%%%%%%%%%%%%%%%%%%%%%%%%%%%%%%%%%%%%%%%%%%%%%%%%%%%%%%%%%%%%%%%%%%%%%%%%%
%%%%%%%%%%%%%%%%%%%%%%%%%%%%%%%%%%%%%%%%%%%%%%%%%%%%%%%%%%%%%%%%%%%%%%%%%%
\subsubsection{Basic properties}


The stochastic perturbations defined in \eqref{eq: definition combined noise} form a martingale difference sequence with bounded second moment.

\begin{lemma}
% [Unbiased perturbations with bounded conditional variance]
\label{thm: properties of noise}
% Under initial- or first-visit updates, 
There exists a constant $c_w \in [0, \infty)$ 
% independent of $\mu_k$
such that for every time $k \ge 0$, every state $s \in \sset$ and every action $a \in \aset(s)$,
\begin{align}
\label{eq: muw_mds}
% \E\!\left[\mu_k(s,a)\,w_k(s,a)\mid \mathcal F_k\right] &= 0,\\
&\mathbb E\!\left[w_k(s,a) \mid \mathcal F_k\right] = 0,
\\
\label{eq: muw_second_moment}
% \E\!\left[\mu_k^2(s,a)\,w_k^2(s,a)\mid \mathcal F_k\right]
&\mathbb E\!\left[w_k^2(s,a) \mid \mathcal F_k\right]
\le c_w(1+\|q_k\|^2) .
\end{align}
% In particular, $\mu_k(s,a)w_k(s,a)$ has bounded conditional variance uniformly in $k$ and $(s,a)$.
\end{lemma}


Using \autoref{thm: properties of noise} and the Robbins--Monro conditions,
we can apply standard martingale and stochastic approximation arguments to establish fundamental properties of MC-O-PI.
% For simplicity, we state these results under \autoref{asp: standing}, even though they hold under weaker conditions. We defer the discussion of these relaxations to Section \ref{sec: relaxing conditions} and the proofs to \autoref{sec: Appendix}.


% \begin{lemma}
% \label{eq: properties of noise} \textcolor{red}{C(1+\|q_k\|)}
%     Under \autoref{asp: standing}
%     there exist two constants $c_1, c_2 \in \real_{\geq 0}$ independent of $\update$
%     such that
%     for every point $q \in \qset$, policy $\pol \in \polset$, state $s \in \sset$ and action $a \in \aset(s)$,
%     the noise $w$ defined in equation \eqref{eq: definition combined noise} satisfies
%     \begin{align}
%         \label{eq: noise prop exp}
%         &\E\big[ \update(s,a) w(s,a) \mid \mathcal{F}_k \big] = 0 ,
%         \\
%         \label{eq: noise prop var}
%         &\E\big[ \update^2(s,a) w^2(s,a) \mid \mathcal{F}_k \big] \leq c_1 + c_2 \| q \|^2 . 
% \end{align}
% \end{lemma}
% \begin{proof}
% Since the variables $v$ and $v'$ in equation \eqref{eq: definition combined noise} are independent, we know that
% \begin{align*}
%     \E \big[ v v' \big]
%     = \E \big[ v \big]  \E \big[ v' \big] = \zerofun 
% \end{align*}
% where $\zerofun$ is the constant function $\zerofun(s,a) = 0$.
% Thus it follows that
% \begin{align*}
%     \E \big[ \update w \big] 
%     &= \E \big[ \update v + v' ( q_\pol - q + v ) \big] 
%     \\
%     &= \update \E[v] + \E [v'] \big( q_\pol - q + \E[ v ] \big) 
%     \\
%     &= \zerofun.
% \end{align*}
% Further, we know that for all state-action pairs $0 \leq \update(s,a) \leq 1$ and $-1 \leq v'(s,a) \leq 1$. Thus we obtain that
% \begin{align*}
%     \E \big[ \update^2 w^2 \big]
%     &= \E \big[ ( \update v + v' ( q_\pol - q + v ) )^2 \big]
%     \\
%     % &= \E \Big[ w_1^2 + \frac{w_2^2}{\update^2} ( q_\pol - q + w_1 )^2 + 2 \frac{w_1 w_2}{\update} ( q_\pol - q + w_1 ) \Big]
%     % \\
%     % &= \E \big[ w_1^2 \big] + \E \Big[ \frac{w_2^2}{\update^2} ( q_\pol - q + w_1 )^2 \Big] + 2 \E \Big[ \frac{w_1 w_2}{\update} ( q_\pol - q ) \Big] + 2 \E \Big[ \frac{w_1^2 w_2}{\update} \Big]
%     % \\
%     &= \update^2 \E \big[ v^2 \big] + \E \big[ v'^2 ( q_\pol - q + v )^2 \big] +  2 \E \big[ \update v v' ( q_\pol - q + v ) \big] 
%     \\
%     &= \update^2 \E \big[ v^2 \big] + \E \big[ v'^2 ( q_\pol - q + v )^2 \big]
%     \\
%     &\leq \E \big[ v^2 \big] + \E \big[ ( q_\pol - q + v )^2 \big]
%     \\
%     &= \E \big[ v^2 \big] + (q_\pol - q)^2 + \E \big[ v^2 \big] + 2 (q_\pol - q) \E \big[ v \big]
%     \\
%     &\leq c_1 \onefun + c_2 q^2
%     % \\
%     % &\leq (c_1 + c_2 \|q\|^2) \onefun
% \end{align*}
% for some constants $c_1, c_2 \in \real_{\geq 0}$ when every action-value function $q_\pol$ is bounded.
% Finally, \eqref{eq: noise prop var} follows from the inequality $q(s,a)^2 \leq \|q\|^2$. and chose $C = \max[c_1, c_2]$.
% \end{proof}



\begin{lemma}
% [Bounded solutions]
\label{thm: bounded limit sets}
Let $\polset_\infty$ denote the set of policies that appear infinitely often in a sequence $(\pi_k)_{k \ge 0}$ generated by MC-O-PI.
% \begin{align}
%     \polset_\infty \coloneqq \bigcap_{t = 0}^{\infty} \text{cl} \{\pi_k : k \geq t \} ,
% \end{align}
% Under \autoref{asp: standing},
The corresponding sequence $(q_k)_{k \ge 0}$ converges almost surely to the set
\begin{align*}
    \Big\{
    q \in \qset : 
    \min_{\pi \in \polset_\infty} q_\pi(s,a) \leq q(s,a) \leq \max_{\pi \in \polset_\infty} q_\pi(s,a), \forall s \in \sset, \forall a \in \aset
    \Big\} .
\end{align*}
Thus
there exists an almost surely finite constant $c_Q \in [0, \infty)$ such that for all $k \geq 0$
    \begin{align*}
        \| q_k \|_\infty^2 < c_Q .
    \end{align*}
\end{lemma}

% \begin{corollary}
% Under \autoref{asp: standing},
% The limit set of any solution is contained in 
% \begin{align}
%     \{
%     q \in \qset : 
%     q_{\worst{\pi}} \leq q \leq \qopt
%     \}
% \end{align}
% \end{corollary}

% Given that there are a finite number of deterministic policies, at least one of them must be visited infinitely often.
% As a result, \autoref{thm: bounded limit sets} implies that solutions are almost surely bounded.

% \begin{corollary}[Solutions are bounded a.s.]
% \label{thm: bounded solutions}
% For every solution $(q_k)$ of MC-O-PI, there exists under \autoref{asp: standing} almost surely a random constant $M \in \real_{\geq 0}$ such that for all $k \geq 0$
%     \begin{align}
%         \| q_k \|_\infty^2 < M .
%     \end{align}
% \end{corollary}


% Given that MC-O-PI solutions are bounded, the variance of the stochastic perturbation in \eqref{eq: muw_second_moment} remains bounded as well.
% As a result, the total effect of the perturbations is finite.


% % \textcolor{red}{the effects of the noise vanish accumulation
% % the effect of the noise dont vanish relative to the update, because in order to paply Teel framework would need perturbations to scale down relative to the update}

% \begin{lemma}
% [Bounded total perturbations]
% \label{thm: effects of noise decrease}
% \textcolor{red}{only keep lemma if necessary for explaining relaxed conditions in discussion}
% Under \autoref{asp: standing},
% for every initial condition $q_0 \in \qset$, 
%     the series $\sum_{k=0}^\infty \lrate_k w_k$ generated by MC-O-PI converges almost surely to a finite limit.
%     % and update distributions $(\update_k)_{k \geq 0}$ satisfy the modified Robbins--Monro conditions (\autoref{asp: stochastic conditions on effective learning rate}).
%     Hence the sequence $(\lrate_k w_k)_{k \ge 0}$ converges almost surely to zero.
% \end{lemma}


% As solutions converge to a bounded region 
% and the effects of noise become arbitrarily small, 
% updates of the differential inclusion in \eqref{eq: mcopi sa} become arbitrarily small. 
% In other words, the distance between consecutive elements of a solution tends to zero.

% \begin{lemma}
% \label{thm: steps become small}
% For every solution $(q_k)_{k \ge 0}$ of MC-O-PI,
% % Let $(q_k)_{k \geq 0}$ be a sequence generated by MC-O-PI, as defined in \eqref{eq: mcopi sa}.
% the sequence $( \|q_\kpo - q_k\| )_{k \geq 0}$ converges almost surely to zero under \autoref{asp: standing}.
% \end{lemma}


%%%%%%%%%%%%%%%%%%%%%%%%%%%%%%%%%%%%%%%%%%%%%%%%%%%%%%%%%%%
%%%%%%%%%%%%%%%%%%%%%%%%%%%%%%%%%%%%%%%%%%%%%%%%%%%%%%%%%%%
% \subsubsection{Policy transition}

If the target action values $q_{\pi_k}$ in \eqref{eq: mcopi sa} were fixed,
solutions would almost surely converge to this fixed target by Proposition~4.1 or Example~4.3 in~\citep{Bertsekas1996Neuro-DynamicProgramming}.
Thus, the asymptotic behaviour of the difference inclusion hinges on how these action values evolve over time.
To isolate these changes, we introduce \textit{transition times}.

% \noindent
% {\bf Theorem} {\it Let $u,v,w$ be discrete variables such that $v, w$ do
% not co-occur with $u$ (i.e., $u\neq0\;\Rightarrow \;v=w=0$ in a given
% dataset $\dataset$). Let $N_{v0},N_{w0}$ be the number of data points for
% which $v=0, w=0$ respectively, and let $I_{uv},I_{uw}$ be the
% respective empirical mutual information values based on the sample
% $\dataset$. Then
% \[
% 	N_{v0} \;>\; N_{w0}\;\;\Rightarrow\;\;I_{uv} \;\leq\;I_{uw}
% \]
% with equality only if $u$ is identically 0.} 
% \hfill\BlackBox


\begin{definition}
\label{def: transition times}
Given a sequence of policies $(\pi_k)_{k \ge 0}$,
define the transition times $(t_n)_{n \ge 0}$ as
\begin{align*}
t_0 &:= 0 ,
\\
t_{n+1} &:= \inf\{k>t_n : q_{\pi_k} \neq q_{\pi_{t_n}} \}.
\end{align*}
Thus, if $t_n = \infty$, then $t_{n+i} = \infty$ for all $i \ge 1$.
These times decompose the sequence $(\pi_k)_{k \ge 0}$ into blocks of constant action values, such that $q_{\pi_k} = q_{\pi_{t_n}}$ for all $k \in \{t_n, \dots, t_{n+1}-1 \}$.
\end{definition}

% \begin{lemma}
% Let \((\pi_k)_{k\ge 0}\) be a sequence of policies generated by MC-O-PI,
% and let \((t_n)_{n\ge 0}\) be the associated transition times.
% Under \autoref{asp: standing},
% if for some index $n \in \mathbb Z_{\ge 0}$, the policy $\pi_{t_n}$ is suboptimal, then
% \[
% \mathbb P( t_{n+1}<\infty \mid \mathcal F_{t_n} ) = 1.
% \]
% \end{lemma}
% \begin{proof}
% Let \((\pi_k)_{k\ge 0}\) be a sequence of policies generated by MC-O-PI,
% and let \((t_n)_{n\ge 0}\) be the associated transition times.

% If \(t_{n+1}=\infty\), then by \autoref{def: transition times},
% \(q_{\pi_k}=q_{\pi_{t_n}}\) for all \(k\ge t_n\).
% Hence the target remains fixed forever, so \(q_k\to q_{\pi_{t_n}}\) almost surely.

% Since \(\pi_{t_n}\) is suboptimal, it is not greedy with respect to \(q_{\pi_{t_n}}\). Thus there exist
% \(s\in S\) and \(a\in A(s)\) such that
% \[
% q_{\pi_{t_n}}(s,a) > q_{\pi_{t_n}}(s,\pi_{t_n}(s)).
% \]
% By convergence, the same strict inequality holds for all large \(k\), so
% \(\pi_k(s)\neq \pi_{t_n}(s)\) for all large \(k\) because \(\pi_k\) is greedy
% with respect to \(q_k\). Therefore \(q_{\pi_k}\neq q_{\pi_{t_n}}\) for some
% \(k>t_n\), contradicting \(t_{n+1}=\infty\). Hence \(t_{n+1}<\infty\) almost
% surely.
% \end{proof}



\begin{lemma}
\label{thm: suboptimal blocks finite}
If $\pi_{t_n}$ is not optimal, then $t_{n+1}<\infty$ almost surely.
\end{lemma}
\begin{proof}
Let $\pi := \pi_{t_n}$.
If $t_{n+1}=\infty$, then the target is fixed at $q_\pi$ forever and $q_k \to q_\pi$ almost surely~\cite[Proposition~4.1]{Bertsekas1996Neuro-DynamicProgramming}.
Since $\pi$ is suboptimal, it is not greedy with respect to $q_\pi$.
Hence there exists a pair $(s,a)$ such that $q_\pi(s,a) > q_\pi(s,\pi(s))$.
Consequently, as $q_k$ approaches $q_\pi$,
every greedy policy $\pi_k$ with respect to $q_k$ satisfies
$q_\pi(s,\pi_k(s)) > q_\pi(s,\pi(s))$
for all sufficiently large $k$.
The Policy Improvement Theorem indicates that $q_{\pi_k} \ne q_\pi$, contradicting $t_{n+1}=\infty$.
\end{proof}

% Clearly, if $\pi_{t_n}$ is suboptimal, then $t_{n+1}<\infty$.
% Otherwise the target would remain equal to \(q_{\pi_{t_n}}\) forever, forcing convergence to \(q_{\pi_{t_n}}\).
% However a suboptimal policy is not greedy with respect to its own action values.
% Thus no solution can get arbitrarily close to $q_{\pi_{t_n}}$ without causing the greedy policy to improve. Therefore no suboptimal block can persist indefinitely.
% 
\autoref{thm: suboptimal blocks finite} implies that if there exists a finite index $n$ such that $t_{n} = \infty$, the estimate $q_k$ converges to the action values of the last bounded transition time, which must be optimal.
Conversely, any suboptimal asymptotic behaviour requires an infinite sequence of bounded transition times.
% 
% These conclusions remain valid in the stochastic setting of \eqref{eq: mcopi sa} since the noise is a martingale difference sequence \cite[Proposition~4.1]{Bertsekas1996Neuro-DynamicProgramming}.

% We rely on the following result at the transition times.

% \begin{lemma}
% \label{thm: strict improvement theorem}
% For all policies $\pi, \pi' \in \polset$ that satisfy
% \begin{align}
% \label{eq: strict improvement theorem}
% q_\pi(s,\pi) \le q_\pi(s,\pi')
% \end{align}
% for every state $s \in \sset$,
% if $q_\pi \neq q_{\pi'}$, then policy $\pi'$ is better than $\pi$ because $v_{\pi} \le v_{\pi'}$ and $v_{\pi}(s) < v_{\pi'}(s)$ for at least one state $s \in \sset$.
% \end{lemma}



%%%%%%%%%%%%%%%%%%%%%%%%%%%%%%%%%%%%%%%%%%%%%%%%%%%%%%%%%%%%%%%%%%%%%%%%%%
%%%%%%%%%%%%%%%%%%%%%%%%%%%%%%%%%%%%%%%%%%%%%%%%%%%%%%%%%%%%%%%%%%%%%%%%%%
% \subsection{Convergence to optimality}

% Example~5.12 in \citep{Bertsekas1996Neuro-DynamicProgramming} shows that the Robbins--Monro conditions alone do not prevent initial-visit MC-O-PI from converging to a suboptimal limit point.
% More broadly, it is straightforward to introduce suboptimal limit points, even on the boundary of the optimal region $\graval(\pi_*)$, for instance, by strongly biassing the update distributions towards the non-greedy actions at every step.

% \textcolor{red}{Recent counterexamples show that the exploring-starts condition alone is not sufficient ... (cite us)?}

% % Global convergence therefore requires additional assumptions, either on the MDP structure or on the sampling and update schemes of MC-O-PI.

% % For instance, 
% % \citet{Wang2022OnLearning} prove convergence to the optimal action values provided no state is visited twice when simulating an episode under the optimal policy $\best{\pol}$. Yet this feedforward condition is unverifiable when the MDP is unknown.

% % % also examples
% % % where
% % % \cite{Harutyunyan2016QCorrections} proof for every visit with discount < 1/3
% % % \cite{Delattre2025MarkovAlgorithm} proof for first visit with discount < 1/2
 
% % Alternatively, \citet{Tsitsiklis2002OnIteration} and later \citet{Chen2018OnProblem} and \citet{Liu2021OnStarts} prove convergence to the optimal action values provided all state-action pairs are updated equally often on average. In other words, they require that the update distributions be uniform over all state-action pairs.

% \citet{Tsitsiklis2002OnIteration} and later \citet{Chen2018OnProblem} and \citet{Liu2021OnStarts} prove convergence to the optimal action values provided all state-action pairs are updated equally often on average. In other words, they require that the update distributions be uniform over all state-action pairs.

% This paper extends theirs results.

% The proofs \cite{Tsitsiklis2002OnIteration}

% where $T$ and $P_\pi$ denote the Bellman optimality operator and the transition operator under policy $\pi$ ()

% The key insight of \cite{Tsitsiklis2002OnIteration} is that 
% when $\update_k$ is proportional to the identity map, it commutes with the transition operator $P_{\pi_k}$.
% As a result, MC-O-PI solutions satisfy
% \begin{align}
% \label{eq: q-Bq decreases}
%     q_\kpo - T q_\kpo \leq (I- \lrate_k \update_k)(q_k - T q_k) - \lrate_k (\discount P_{\pol_k} \update_k w_k -  \update_k w_k)
% \end{align}

% The Robbins--Monro learning rates and the martingale difference perturbations then ensure that the Bellman residual is asymptotically bounded above by zero for every state-action pair:
% \begin{align}
% \label{eq: tsitsiklis central result}
%     \limsup_{k \to \infty} \big( q_k(s,a) - (T q_k)(s,a) \big) \leq 0 , 
%     \qquad \forall s \in \sset, \forall a \in \aset(s).
% \end{align}
% This asymptotic bound implies that for sufficiently large $k$, we have $q_k \le T q_k$.
% By the standard contraction and monotonicity properties of the Bellman operators, we find
% \begin{align*}
%     q_k \leq T q_k = T_{\pol_k} q_k \leq T_{\pol_k}^2 q_k \leq \dots \leq q_{\pol_k} .
% \end{align*}
% Thus when all state-action pairs are updated equally often on average, MC-O-PI solutions converge to the optimal action values because the estimates $(q_k)$ are eventually consistently updated upwards while remaining bounded above.

